%!TEX root = ../sztuthesis_main.tex
\phantomsection
\addcontentsline{toc}{section}{Abstract}

\keywordsen{music emotion analysis; large language model; audio description generation; Valence-Arousal; DEAM dataset}
\categoryen{TP391}
\begin{abstracten}
With the development of music content analysis and intelligent recommendation technologies, how to effectively extract semantic information from music audio and further perform emotion recognition has become an important issue in music information retrieval and affective computing. Traditional music emotion analysis methods mainly rely on acoustic features and machine learning models. Although they are relatively stable in numerical prediction, they are limited in high-level semantic representation and interpretability. To address this problem, this thesis studies music audio description and emotion label generation based on large language models, and proposes a two-stage processing framework. In the first stage, a pretrained audio understanding model is used to generate Chinese textual descriptions for music audio clips. In the second stage, a large language model performs emotion inference based on the generated descriptions and outputs Valence, Arousal, and discrete emotion labels.

This work takes the DEAM dataset as the experimental basis and completes audio preprocessing, sample construction, and the division of training, validation, and test sets. In the audio description generation stage, music clips are cropped, resampled, and encoded, and prompt constraints are designed to guide the model to generate descriptions from five dimensions: melody, rhythm, instruments, style, and overall listening impression. In addition, a rule-based quality assessment method is designed to automatically evaluate the generated descriptions in terms of length, dimensional coverage, and hallucination risk. In the emotion inference stage, music emotion analysis is modeled as a semantic reasoning task based on textual descriptions. A large language model is employed to output continuous emotion dimensions and discrete emotion labels, while engineering functions such as batch inference, result parsing, and resume-from-breakpoint are also implemented. To verify the feasibility of the proposed framework, a traditional machine-learning baseline based on acoustic features is further constructed for comparison.

Experimental results show that the proposed system is able to complete the full processing pipeline from music audio input to textual description generation and then to emotion label output. On a small-scale test set, the generated descriptions are generally usable and can provide semantic support for subsequent emotion inference. Comparative experiments indicate that the traditional baseline is still more stable in numerical prediction at the current stage, while the proposed large-model-based method has advantages in semantic representation and interpretability of the intermediate process. Overall, this thesis verifies the feasibility of combining an audio-description intermediate layer with large language model based emotion inference for music emotion analysis, while also revealing that further improvements are needed in experimental scale, evaluation completeness, and performance optimization.
\end{abstracten}
